\documentclass{report}
\usepackage[top=1in, bottom=1in, left=1in, right=1in]{geometry}
\usepackage{graphicx}
\usepackage{amsmath}
\usepackage{amssymb}
\usepackage{amsthm}
\usepackage{algorithm}
\usepackage{algpseudocode}
\usepackage{hyperref}
\hypersetup{colorlinks=true, linkcolor=blue, urlcolor=blue, citecolor=blue}
\usepackage{enumitem}
\usepackage{tikz}
\usetikzlibrary{positioning, arrows.meta}


\newtheorem{claim}{Claim}

\title{Communication Efficient Gradient Coding: A Survey}
\author{Jay Mehta: 22B1281\\
        Kshitij Vaidya: 22B1829\\
        Tanay Bhat: 22B3303}
\date{\today}  

\newtheorem{theorem}{Theorem}
\newtheorem{lemma}{Lemma}

\renewcommand\qedsymbol{$\blacksquare$}
\begin{document}

\maketitle

\tableofcontents
\chapter{Implementation}
In this chapter, we will discuss our own implementation of a toy example of a gradient coding scheme. We will implement a simple gradient coding scheme based on the idea of polynomial interpolation as described in the reference paper. We will use Python for our implementation.
\section{Setup}
We have made an assumption in our implementation.
\begin{equation*}
    \text{n (Number of workers)} = \text{k (Number of datasets)}
\end{equation*}
This is consistent with the reference paper, which also makes a similar assumption for simplicity. 

We test the correctness of our implementation by trying out multiple test cases mentioned below:
\begin{table}[!h]
    \centering
\begin{tabular}{c|c|c|c|c|c}
\hline
\textbf{Test Case} & \textbf{n = k} & \textbf{d} & \textbf{s} & \textbf{m} & \textbf{l} \\ \hline
1 & 5 & 3 & 1 & 2 & 4 \\
2 & 6 & 3 & 1 & 2 & 6 \\
3 & 8 & 4 & 2 & 2 & 8 \\
4 & 10 & 5 & 3 & 2 & 20 \\
5 & 10 & 4 & 2 & 2 & 20 \\
6 & 10 & 3 & 1 & 2 & 20 \\
7 & 12 & 6 & 3 & 3 & 12 \\
8 & 15 & 5 & 2 & 3 & 15 \\ 
\hline
\end{tabular}
\caption{Test cases for our implementation}
\end{table}

n = number of workers, k = number of datasets, d = number of datasets assigned to each worker, s = number of stragglers tolerated, m = communication reduction factor, l = dimension of the gradient vector.
\end{document}